\begin{trapbox}
\begin{enumerate}
\item \textbf{把顶角当锐角。}
$A_1,A_2$ 是长轴端点时，$P$ 在非圆椭圆上非端点，
$\angle A_1PA_2$ \textbf{始终钝角}，不是锐角。
因此 $\tan\alpha<0$。

\item \textbf{最大角公式丢掉正切的负号。}
正确式为
$\tan\alpha_{\max}=-2ab/(a^2-b^2)$。
正数 $2ab/(a^2-b^2)$ 是\textbf{补角的正切值}：
$\tan(180^\circ-\alpha_{\max})$。

\item \textbf{认为 $\arctan(4/3)$ 是最大顶角。}
$\arctan(4/3)$ 是第一象限锐角，
真正最大顶角为 $180^\circ-\arctan(4/3)$，
也等于 $2\arctan2$（$a=2,b=1$ 时）。

\item \textbf{把 $\tan\alpha$ 的绝对值当作随高度增大而增大。}
事实上
$|\tan\alpha|=2ab^2/(c^2|y|)$，
它随 $|y|$ 增大而\textbf{减小}；
但 $\alpha$ 随 $|y|$ 增大而\textbf{增大}。
最好改用 $\tan(\alpha-90^\circ)$ 的正向单调关系。

\item \textbf{把两条直线的锐夹角当三角形顶角。}
两条相交直线有一对锐角、一对钝角。
本题要求的是 $P$ 处两条\textbf{射线}形成的普通内角，
需要用点积、几何位置或底角和辨认象限。

\item \textbf{把长轴端点写成“最小值取得处”。}
$P=A_1$ 时射线 $PA_1$ 长度为零，$\angle A_1PA_2$ 无定义；
$P=A_2$ 也一样。严格说
$\inf\alpha=90^\circ$，\textbf{不取到}。

\item \textbf{斜率积恒定不等于夹角恒定。}
虽有 $k_1k_2=-b^2/a^2$，但
$\tan$ 两直线夹角还依赖于 $k_1-k_2$。
只在圆的特殊情形 $a=b$ 才因直径圆周角使顶角恒为直角。

\item \textbf{已知角度就认为点唯一。}
对于固定椭圆的 $90^\circ<\alpha<\alpha_{\max}$，
通常对应四个点 $(\pm|x|,\pm|y|)$；
最大角只对应两个短轴端点。

\item \textbf{拿焦角公式套顶角。}
C003 用的是焦点 $F_1,F_2$、底边 $2c$，面积是 $c|y|$；
本题底边 $A_1A_2=2a$，面积是 $a|y|$。
两种半角关系与角度范围不可混用。

\item \textbf{圆的退化情形忽略分母零。}
当 $a=b$ 时，顶角恒为 $90^\circ$，正切无定义。
不可向 $\tan\alpha=-2ab^2/((a^2-b^2)|y|)$ 直接代 $a=b$。

\item \textbf{认为“离长轴越远，视角越大”对平面任意点成立。}
这种严格单调性是\textbf{$P$ 被约束在同一个椭圆上}的结论。
如果 $P$ 任意变化横坐标，纵坐标大小不能单独决定视角。

\item \textbf{角度存在性只解方程，不先验范围。}
给 $60^\circ$、$90^\circ$：本题\textbf{根本无解}；
给 $\alpha_0>90^\circ$：先查
$\alpha_0\le2\arctan(a/b)$，再求坐标。

\item \textbf{直接把 $\tan90^\circ$ 记成 $+\infty$ 并代入差角公式。}
$\tan90^\circ$\textbf{无定义}，不能代入。
可改用正弦、余弦差角公式；若使用极限法，
必须写清 $\beta\to90^{\circ-}$，再令
$t=\tan\beta\to+\infty$，同除以 $t$ 后求极限。

\item \textbf{错记 $\tan(\alpha-90^\circ)=\cot\alpha$ 或把减号当图像左移。}
正确为 $\tan(\alpha-90^\circ)=-\cot\alpha$。
本题 $\alpha$ 为钝角，左边是锐角正切（正数），
而 $\cot\alpha$ 是负数，\textbf{用符号即可排错}。
函数自变量写 $x-90^\circ$，图像对应\textbf{向右}平移 $90^\circ$。
\end{enumerate}
\end{trapbox}
